%% CV Mestre - Laura Lacort Zimmermann
%% Visão abrangente de todas as competências, experiências e conquistas.
%% Usar como referência-mestre ao adaptar CVs por vaga.
%%
%% Compile com: lualatex main_example.tex
%% (pdflatex frequentemente falha em instalações MiKTeX modernas com erros de
%% expansão de fonte do fontawesome5.)

\documentclass[11pt,a4paper,sans]{moderncv}
\moderncvstyle{banking}
\moderncvcolor{blue}

% Força nome e sobrenome, além dos títulos de seção, a renderizar na cor azul
% do moderncv (color1). Por padrão, em lualatex+MiKTeX, ficam pretos, o que
% destoa do restante do esquema de cores em azul.
\renewcommand*{\firstnamestyle}[1]{{\fontsize{34}{36}\bfseries\upshape\color{color1}#1}}
\renewcommand*{\lastnamestyle}[1]{{\fontsize{34}{36}\bfseries\upshape\color{color1}#1}}
\renewcommand*{\sectionstyle}[1]{{\sectionfont\color{color1}#1}}

\usepackage[utf8]{inputenc}
\usepackage{hyperref}
\hypersetup{
    colorlinks=true,
    linkcolor=blue,
    filecolor=magenta,
    urlcolor=blue,
    pdftitle={Laura Lacort Zimmermann - CV},
    pdfpagemode=FullScreen,
}
\usepackage[scale=0.80]{geometry}
\usepackage{import}

% dados pessoais
\name{Laura}{Lacort Zimmermann}
\address{Brasília, DF, Brasil}{}{}
\email{lauralacortzimermann@gmail.com}
\extrainfo{\href{https://linkedin.com/in/laurazimrn}{LinkedIn}, \href{https://github.com/laurazimrn}{GitHub}}

\begin{document}

\makecvtitle

% ============================================================
%     DECLARAÇÃO DE PERFIL
% ============================================================

\vspace{6pt}
\small{Desenvolvedora com bacharelado em Ciência da Computação e pós-graduação em Engenharia de Software, com experiência prática em desenvolvimento Front-End e desenvolvimento web. Em transição ativa para as áreas de Backend e Dados/IA, combinando experiência profissional em ciclo completo de desenvolvimento com projetos pessoais em Python (FastAPI, Django, LangChain).}

% ============================================================
%     COMPETÊNCIAS PRINCIPAIS
% ============================================================

\section{Competências Principais}
\vspace{1pt}
\begin{itemize}

\item \textbf{Backend Python}: FastAPI, Django, SQLAlchemy, API REST

\item \textbf{Banco de Dados}: PostgreSQL, MySQL, SQLite, SQL, modelagem relacional

\item \textbf{IA \& LLM}: LangChain, LLM, Machine Learning, Inteligência Artificial

\item \textbf{Frontend}: JavaScript, TypeScript, HTML5, CSS3, React, Tailwind CSS

\item \textbf{Ferramentas}: Git, GitHub, Jira, Trello, Power BI, Salesforce, Figma, Postman, Insomnia, Linux, Bubble.io

\end{itemize}

% ============================================================
%     EXPERIÊNCIA PROFISSIONAL
% ============================================================

\section{Experiência Profissional}
\vspace{3pt}
\begin{itemize}

\item{\cventry{2024--Atual}{Desenvolvedora No Code}{IMMA Experiências}{Brasil}{}{\vspace{1pt}
\begin{itemize}
    \item Desenvolvimento de aplicações web completas utilizando Bubble.io, com foco em modelagem de banco de dados relacional e definição de estruturas de dados
    \item Aplicação de regras de negócio e integridade referencial no design do banco de dados
    \item Integração com APIs REST externas e construção de fluxos de automação
    \item Entrega de interfaces responsivas alinhadas a requisitos funcionais
    \item Atuação em ciclo completo de desenvolvimento: levantamento de requisitos, modelagem, implementação e testes
    \item Uso diário de Jira para acompanhamento de demandas e projetos
\end{itemize}}}

\vspace{3pt}

\item{\cventry{2023}{Desenvolvedora Front End}{Agência Sete Digital}{Brasil}{}{\vspace{1pt}
\begin{itemize}
    \item Desenvolvimento de interfaces web responsivas com HTML, CSS e JavaScript
    \item Consumo de bibliotecas de componentes UI: Material UI, Tailwind CSS e Bootstrap
    \item Depuração e correção de bugs em ambiente de produção
    \item Colaboração com equipe em ambiente ágil e versionamento com Git/GitHub
\end{itemize}}}

\vspace{3pt}

\item{\cventry{2022}{Desenvolvedora Web}{Eficaz Marketing}{Brasil}{}{\vspace{1pt}
\begin{itemize}
    \item Criação e manutenção de componentes reutilizáveis para interfaces web
    \item Implementação de layouts responsivos seguindo padrões de UI/UX
    \item Versionamento de código com Git/GitHub
\end{itemize}}}

\vspace{3pt}

\item{\cventry{2021--2022}{Estágio Front End}{Eficaz Marketing}{Brasil}{}{\vspace{1pt}
\begin{itemize}
    \item Desenvolvimento de interfaces web com HTML, CSS e JavaScript
    \item Manutenção de sites e aplicativos
    \item Versionamento de código com Git/GitHub
\end{itemize}}}

\end{itemize}

% ============================================================
%     PROJETOS PESSOAIS
% ============================================================

\section{Projetos Pessoais}
\vspace{1pt}
\begin{itemize}

\item \textbf{Local AI Agent com Ollama \& LangChain}: agente de IA local em Python com LangChain, Ollama e ChromaDB (RAG). Processa arquivos CSV e realiza raciocínio contextual sem APIs externas.

\item \textbf{Notes API with AI}: API RESTful com FastAPI e SQLAlchemy integrada ao Ollama (llama3) para sumarização de notas. CRUD completo, Pydantic v2 e testes com Pytest.

\item \textbf{FastAPI App with PostgreSQL}: aplicação de quiz com FastAPI, PostgreSQL e SQLAlchemy. Modelagem relacional e documentação interativa de endpoints.

\item \textbf{Django REST API}: API RESTful com Django REST Framework, arquitetura modular. Banco SQLite com suporte a PostgreSQL.

\item Disponíveis em: \href{https://github.com/laurazimrn}{github.com/laurazimrn}

\end{itemize}

% ============================================================
%     EDUCAÇÃO
% ============================================================

\section{Educação}
\vspace{1pt}
\begin{itemize}

\item{\cventry{2026--2027}{Especialização em Matemática, Educação e Tecnologia}{IFB}{Brasil}{}{\vspace{1pt}
Em andamento.
}}

\vspace{3pt}

\item{\cventry{2024--2025}{Pós-graduação em Engenharia de Software}{UTFPR}{Brasil}{}{\vspace{1pt}
Concluído.
}}

\vspace{3pt}

\item{\cventry{2021--2023}{Bacharelado em Ciência da Computação}{UNIVEM}{Brasil}{}{\vspace{1pt}
Concluído.
}}

\end{itemize}

% ============================================================
%     IDIOMAS
% ============================================================

\section{Idiomas}
\vspace{1pt}
\begin{itemize}
\item Português (nativo), Inglês (intermediário), Espanhol (básico).
\end{itemize}

% ============================================================
%     CERTIFICAÇÕES
% ============================================================

\section{Certificações}
\vspace{1pt}
\begin{itemize}
\item{\textbf{AI Fluency Framework \& Foundations} -- Anthropic (mar 2026).}
\item{\textbf{Data Fundamentals} -- IBM (fev 2026).}
\item{\textbf{Python Essentials 1} -- Cisco Networking Academy (dez 2025).}
\item{\textbf{Python Essentials 2, JavaScript Essentials 1, Fundamentos de Data Science e IA} -- em andamento.}
\end{itemize}

% ============================================================
%     REFERÊNCIAS
% ============================================================

\section{Referências}
\vspace{1pt}
\begin{itemize}
\item{Disponíveis mediante solicitação.}
\end{itemize}

\end{document}
